\chapter{Rechtliche Grundlagen -- \gAasRsN{RS~III}} 

\emph{Christof Koller}

\gAuthNotes{
Changelog:
2009-02-20: GABS: Fertig f\"ur Kurs 2009/02

2009-02-25: EDIT Koller

2009-04-22: GABS: Abkürzungen "`Pat."', "`KA"', "`KH"'
bereinigt und gegen Langform ausgetauscht, wie besprochen.}

\setcounter{tocdepth}{10}
\subsection*{Übersicht}
\mtcsetdepth{minitoc}{3}
\minitoc
\mtcsetdepth{minitoc}{2}



\section{Allgemeine Rechtsgrundlagen}
\subsection{Struktur der \"Osterreichischen Rechtsordnung}

%\begin{tabular}{| p{2.2cm} | p{2.2cm} p{2.2cm} p{2.2cm} | p{2.2cm} p{2.2cm}
% |}\toprule

\begin{table}[H]
\gCaptionof{table}{Rechtsordnung}
\begin{tabularx}{\linewidth}[H]{X X|X|X||X|X||}

& \multicolumn{5}{c}{\gEs{$\Swarrow$ EU-Recht $\Searrow$}} 
\\\addlinespace

& \multicolumn{5}{c}{\gEs{Struktur der \"Osterreichischen Rechtsordnung}} 
\\\addlinespace

{ } 
& \multicolumn{3}{c}{\centering \"Offentliches Recht}
& \multicolumn{2}{c}{\centering {Privatrecht} } 
\\\hhline{~-----}

\multicolumn{1}{p{2cm}||}{Rechts\-gebiete }
& Strafrecht
& VwR
& VfR
& Zivilrecht
& Sonder\-privat\-rechte \\\hhline{~-----}

\multicolumn{1}{p{2cm}||}{Wichtigste Gesetzesbereiche} & 
	StGB, StPO &
		Berufs\-ge\-setze,\newline KAKuG, UbG & 
			BVG; StGG &
				ABGB & 
					Arbeits\-rechts\-ge\-setze,
					Unter\-nehmens\-ge\-setz\-buch\\\hhline{~-----}

\end{tabularx}
\end{table}
%\ncbarr[arrows=->]{EU}{OR}
% TODO REcht: Linien: Dicker statt doppelt

Die \"osterreichische Rechtsordnung ist ein zusammenhängendes Geflecht
bestehend aus zwei Hauptbereichen, dem \"offentlichen Recht und dem
Privatrecht. Während das \"offentliche Recht die Beziehung zwischen
Staat und B\"urger regelt, bestimmt das Privatrecht die Verhältnisse
unter den B\"urgern (= Privaten) selber: 

\begin{description}
	\item[\gT{\"Offentliches Recht}:] Staat ${\leftrightarrow}$ B\"urger
	\item[\gT{Privatrecht}:] B\"urger ${\leftrightarrow}$ B\"urger
\end{description}

Durch den Beitritt zur \gE{EU} wurde ein weiterer Gesetzesgeber geschaffen,
der in beide Rechtsgebiete einwirken kann. Unabhängig davon lässt
sich jedes von beiden Rechts\-ge\-bieten in folgende
Bereiche unterteilen:

\paragraph*{\"Offentliches Recht}

\begin{itemize}
	\item \gT{Strafrecht}: Erfasst jene Delikte, die der
	Gesetzgeber als derart schwerwiegend erachtet hat, dass
	sie nicht nur (idR.) vom Staatsanwalt vor Gericht zu
	verfolgen sind,  sondern auch (in den meisten Fällen) mit
	Haftstrafen  sanktionierbar sind. Der Staatsanwalt
	vertritt hierbei  den Staat als Ankläger. Die eingehobene
	Geldstrafe  erhält der Staat, nicht der Verletzte.

	\item \gT{Verwaltungsrecht}: Umfasst Regelungen die sich idR. nur an bestimmte Personenkreise richten und von Verwaltungsbeh\"orden
(nicht von Gerichten) bei Versto{\ss} insofern sanktioniert werden, dass (vorrangig) Geldstrafen verhängt werden.

	\item \gT{Verfassungsrecht}: Jener Teil der Rechtsordnung, der die grundlegenden Bestimmungen \"uber Staatsgewalt, ihre Organe und
ihre Verhältnisse zu den Staatsb\"urgern sowie Grundsätze der Rechtsordnung enthält.

\end{itemize}

\paragraph*{Privatrecht}

\begin{itemize}
 \item \gT{Zivilrecht}: Enthält jene allgemeinen Normen, die
die Rechtsverhältnisse zwischen gleichgestellten B\"urgern regeln.
Bei Verst\"o{\ss}en muss der Verletzte selber vor Gericht sein Recht
einfordern (kein Staatsanwalt), wodurch er vom Schädiger
Ersatzleistungen f\"ur den Schaden fordern kann (keine Strafe).

 \item \gT{Sonderprivatrechte}: eigenständige Rechtsgebiete,
die mittels selbstständigen Normen besondere Rechtsverhältnisse
zwischen den B\"urgern regeln. Zu den wichtigsten zählen das Arbeits-
und das Unternehmensrecht.

\end{itemize}

\paragraph{Zugang zu den Gesetzestexten}~
% POI 28

\includegraphics[width=\linewidth]{./images/JUS/RisHeader.jpg}

Die Gesetzestexte können im Internet vom Rechtsinformationssystem des Bundes
(RIS) unter der URL \url{http://www.ris.bka.gv.at/} abgerufen werden. 


\subsection{Vertragsverhältnis im Rettungswesen}

\begin{figure}[H]
%     \gAuthNotes{KOLC: Graphik gehört rotiert.
%     
%     GABS: ok, dauert aber noch ein bisserl. Wenn Inhalt
%     passt bitte etwas Geduld.}
\begin{center}
\includegraphics[width=8cm]{./images/recht_vertrag.png}
\end{center}
\gCaptionof{table}{Vertragsverhältnisse im Rettungsdienst}
\end{figure}
Das Rechtsverhältnis im Rettungswesen besteht grundsätzlich aus drei
Parteien: dem Patienten, der Rettungsorganisation
(Einrichtung) und dem (Not-)Arzt oder Sanitäter (Personal).
(In einer Krankenanstalt ist das Verhältnis ähnlich, mit dem Unterschied, dass an die Stelle der
Rettungsorganisation die Krankenanstalt und an jene des Sanitäters
der Krankenpfleger tritt.)

Als wichtigstes Verhältnis in dieser Dreiecksbeziehung ist das
\gEs{Verhältnis zwischen Patient und
Rettungsorganisation} zu bezeichnen, da direkter
(Vertrags-)Partner zum Patient niemals der Sanitäter, sondern stets die von ihm vertretene
Organisation ist. Man nimmt nämlich an, dass ein Patient, wenn er eine
{\quotedblbase}Rettung{\textquotedblleft} verständigt, nicht mit dem
einzelnen Einsatzpersonal in rechtlichen Kontakt treten m\"ochte,
sondern mit der Gesamtheit einer Organisation. (Er bestellt ja nicht
dieses eine Fahrzeug mit jenem speziellen Personal.) Das
Rechtsverhältnis kann nun auf zwei verschiedene Arten zustande
kommen:

\begin{itemize}
	\item {\mdseries
	\gT{ Heilbehandlungsvertrag}: Der Patient selber oder sein
	zuständiger Vertreter	w\"unscht behandelt zu werden}

	\item {\mdseries
	\gEs{\gT{Geschäftsf\"uhrung ohne Auftrag} im Notfall}: der
	Patient befindet sich in einem derartigen Notfall, dass er sich nicht mehr
	äu{\ss}ern kann und dringend eine Behandlung ben\"otigt (Bsp.:
	Bewusstlosigkeit)}

\end{itemize}

Aus beiden Verhältnissen schuldet die Rettungsorganisation die
h\"ochstm\"ogliche Sorgfalt, aber niemals einen bestimmten
Behandlungserfolg. Weiters hat sie vorrangig nach dem
Willen des Patienten vorzugehen und wenn dieser fehlt, nach
seinem hypothetischen Willen.

Im \gEs{Verhältnis zwischen Patient und
Personal} fungiert das Personal
blo{\ss} als Gehilfe zur Erf\"ullung der Pflichten seiner Organisation gegen\"uber
dem Patienten, wodurch man ihn auch als
\gT{Erf\"ullungsgehilfen} bezeichnet. Er ist nicht
Vertragspartner des Patienten, muss aber mit der gleichen Sorgfalt, wie es
von der Organisation verlangt wird, arbeiten. Kommt es zu
Schädigungen des Patienten durch Handlungen des
Personals, kann der Geschädigte deliktisch gegen den
Erf\"ullungsgehilfen und vertraglich gegen die Organisation vorgehen. (IdR. wird der
Patient die Organisation verklagen, da ihm hierbei aufgrund
des Vertrages mehrere gesetzliche Erleichterungen zustehen.)

Das \gEs{Verhältnis zwischen Rettungsorganisation
und Personal} lässt sich auf drei verschiedene Arten
begr\"unden:

\begin{itemize}
\item \gE{Arbeitsvertrag}: Der (Not-)Arzt/Sanitäter ist bei
der Organisation als Hauptamtlicher berufsmä{\ss}ig beschäftigt. 

\item \gE{Vereinsmitgliedschaft}: Der (Not-)Arzt/Sanitäter ist
bei der Organisation als Vereinsmitglied und somit als
Ehrenamtlicher/Freiwilliger tätig. 

\item \gE{Zivildienst}: Der Sanitäter leistet seinen
Zivildienst bei der Organisation.

\end{itemize}

Aus diesem Verhältnis heraus m\"ussen manche der Pflichten, welche die
Org\-anisation gegen\-\"uber dem Patienten treffen, auch
durch das Personal er\-f\"ullt werden. Gleich\-zeitig
treffen auch die Org\-anisation und das Personal Rechte und
Pflicht im Um\-gang mit\-ein\-ander.

\subsection{Strafbarkeitsprüfung}

\begin{enumerate}
\item Tatbestandserfüllung bzw. Schaden
\item Kausalit\"at
\item Rechtswidrigkeit

\begin{itemize}
\item {\mdseries
Rechtfertigungsgr\"unde:}


\begin{itemize}
\item {\mdseries
 Heilbehandlung }

\item {\mdseries
Einwilligung}

\item {\mdseries
Notwehr}

\item {\mdseries
rechtfertigender Notstand}

\item {\mdseries
Anhalterecht Privater}

\end{itemize}
\end{itemize}

\item Verschulden

\begin{itemize}
\item{\mdseries
Entschuldigungsgr\"unde:}


\begin{itemize}
\item Notwehr\"uberschreitung im Affekt
\item entschuldigender Notstand
\item (Unzurechenbarkeit)

\end{itemize}
\end{itemize}
\end{enumerate}

Damit es zur Haftung f\"ur eine gesetzte Handlung kommt, m\"ussen die
oben \"uberblicksartig dargestellten und folgend näher beschriebenen
Voraussetzungen vollständig erf\"ullt sein. Mangelt es auch nur an
einer von ihnen, kommt es zu keiner Haftung. Am häufigsten kommen die
oben aufgezählten Rechtfertigungs- oder Entschuldigungsgr\"unde ins
Spiel.

\paragraph{Tatbestandserf\"ullung }

Erf\"ullung der rechtlichen Beschreibung einer Lebenssituation auf die
eine Rechtsfolge/Strafe angeordnet wird

\paragraph*{Kausalität}

{\mdseries
Die gesetzte Handlung war der Grund f\"ur den Eintritt der
Tatbestandserf\"ullung}

\paragraph*{Rechtswidrigkeit}

Die gesetzte Handlung war objektiv falsch ${\rightarrow}$ Urteil \"uber
Tat ${\rightarrow}$ Versto{\ss} gegen gesetzliche Verbote bzw. Gebote,
vertraglichen Bestimmungen oder die guten Sitten 


\begin{description}
    \item[Rechtfertigungsgr\"unde]  Der begangenen Tat kann kein Urteil gemacht werden, da sich ein
ma{\ss}gerechter Mensch in  derselben Situation gleich verhalten
hätte
\end{description}


\paragraph*{Verschulden}

Die gesetzte Handlung war subjektiv falsch ${\rightarrow}$ Urteil \"uber
Täter

\begin{description}
    \item[Entschuldigungsgr\"unde] \"Uber den Täter kann kein Urteil gemacht werden, da er sich nach
seinen pers\"onlichen   Fähigkeiten und Kenntnissen in dieser
Situation nicht anders verhalten hätte k\"onnen
\end{description}



Die einzelnen Rechtfertigungs- und Entschuldigungsgr\"unde werden in der
Ausarbeitung {\quotedblbase}Ein langer Dienst{\textquotedblleft}
separat behandelt.



\clearpage

\section{Ein langer Dienst}
{\centering
\gEs{{\sffamily\Large Eine Geschichte aus 17 Rechtsfällen}}
\par}

\paragraph*{Zweck}

Behandlung rechtlicher Fragen anhand von Praxisbeispielen, die in einer
zusammen-hängenden Geschichte erzählt werden. Dabei wird
insbesondere versucht auf die Besonderheiten des Rettungswesens Bezug
zu nehmen und die sich daraus ergebenden rechtlichen Probleme einfach
darzustellen. Die Darstellung beinhaltet aber nur einen \"Uberblick
\"uber die behandelten Rechtsgebiete sowie eine Konkretisierung auf die
am häufigsten gestellten Rechtsfragen. 

\paragraph*{Aufbau}

Jeder einzelne Rechtsfall wird in der Geschichte nach folgendem Aufbau
gestaltet:

\begin{itemize}
\item 1) Sachverhalt (= \gE{SV}): Schilderung der Geschichte

\item 2) Rechtsfrage (= \gE{RFr}): Erkennung des sich stellenden Problems

\item 3) Rechtslage (= \gE{RL}): rechtliche Beurteilung des Problems

\item 4) Rechtsfolgen (= \gE{RF}): drohende Konsequenzen und Ausweichm\"oglichkeiten

\end{itemize}

\paragraph*{Vorkommende Personen}

\begin{center}
\tablehead{}\begin{supertabular}{m{6.649cm}m{6.651cm}}
\textbf{\textcolor{red}{A}}lex, Fahrer

\textbf{\textcolor{red}{S}}epp, Transportf\"uhrer

\textbf{\textcolor{red}{B}}ernd, Zivildiener

\textbf{\textcolor{red}{NA}}tascha, Notärztin

 &
Hr. \gE{Hauzu}, 1. Patient

Fr. \gE{Uninformiert}, 2. Patientin

Fr. \gE{Finito}, 3. Patientin

Hr. \gE{Schmerz}, 4. Patient

Fr. \gE{Aufgeregt}, Angeh\"orige

Dr. \gE{Wasbringst}, Aufnahmearzt

Fr. \gE{Hilflos}, 5. Patientin

\\
\end{supertabular}
\end{center}

\paragraph{Der Tag beginnt ...}
Unsere Geschichte beginnt an einem warmen Fr\"uhjahrs\-morgen auf den die
drei Sanitäter Alex, Sepp und Bernd schon lange gehofft hatten, da
sie endlich nach dem langen Winter einen angenehm warmen Sonnentag
genie{\ss}en wollen. Doch wie immer an einem solchen Freudentag, ruft
die Arbeit mit einem ganz normalen Sanitätsdienst und seinen
typischen Problemen.

\subsection{Gefahrenbekämpfung}

\subsubsection{Fall 1. Randalierender Patient}

\paragraph*{Sachverhalt}
Die erste Ausfahrt beginnt nat\"urlich direkt nach Dienstbeginn, ohne
dass wenigstens ein Kaffee getrunken werden hätte k\"onnen. Der auf
unsere RTW-Mann\-schaft wartende Patient, Hr. Hauzu, hingegen hat die
ganze Nacht durchgemacht und ist schon weit\-läufig bekannt f\"ur seine
tobend aggressiven Wut\-aus\-br\"uche im Zuge seiner psychiatrischen
Erkrankung. Als Alex, Sepp und Bernd eintreffen, schmei{\ss}t er mit
mehreren Gegenständen um sich und droht sich oder einen anderen zu
verletzen, falls man sich ihm auch nur nähern w\"urde.

\paragraph*{Rechtsfrage}
Darf man Hr. Hauzu festhalten? Muss Hr. Hauzu eingewiesen werden? 

\paragraph*{Rechtslage}

\gEs{Eine Festhaltung \underbar{darf} erfolgen nach
folgenden Gesichtspunkten}:

\begin{itemize}
\item Notwehr \gZ{\S 3 Abs.~1 StGB}: Abwehr eines Angriffes
auf ein eigenes oder fremdes Rechtsgut, wenn:


\begin{itemize}
\item gegenwärtig oder unmittelbar drohender Angriff,
\item rechtswidrig,
\item von Menschen gesetzt,
\item auf ein notwehrfähiges Gut (= Leben, Gesundheit, Freiheit oder
Verm\"ogen) und
\item das Gelindeste der zur verlässlichen Abwehr verf\"ugbaren Mittel
verwendet wird.

\end{itemize}
\end{itemize}
\begin{itemize}
\item Notwehr\"uberschreitung im Affekt \gZ{{\S}~3 Abs. 2 StGB}:
\begin{itemize}
\item {\mdseries
Sollten Alex, Sepp und Bernd fälschlich einen Angriff angenommen haben,
sind sie entschuldigt, wenn sie aus Furcht gehandelt haben und eine
ma{\ss}gerechte Person sich gleich verhalten hätte. }

\end{itemize}
\end{itemize}

\begin{itemize}
\item  Anhalterecht Privater \gZ{\S 86 Abs. 2 StPO}:
Festhaltung bei Verdacht auf eine gerichtlich strafbare Handlung, welche im vorliegenden Fall
die Gefahr einer K\"orperverletzung (KV) iSd. {\S}{\S} 83 ff StGB und
die Sachbeschädigung iSd. {\S} 125 StGB darstellt. Dabei darf nur auf
angemessene Weise (N\"otigung und leichte KV) angehalten werden, wenn
unverz\"uglich f\"ur eine Meldung an die Polizei gesorgt wird.

\end{itemize}

\gEs{Eine Unterbringung \underbar{muss} erfolgen, nach
folgenden Gesichtspunkten}:

\begin{itemize}
\item \gZ{{\S} 3}~Unterbringungsgesetz (UbG):

\begin{enumerate}
\item psychiatrische Krankheit (= KH) 
\item Selbst- und/oder Fremdgefährdung
\item Verhältnismä{\ss}igkeitsgrundsatz

\end{enumerate}
\end{itemize}
\begin{itemize}
\item zu der Feststellung von a) und b) ben\"otigt es das Gutachten eines
\gE{Amtsarztes}, welcher der Bezirksverwaltungsstelle (in Wien die MA 15)
unterstellt ist und \"uber die Polizei zu verständigen ist. Der
Amtsarzt darf die {\quotedblbase}Zuweisung{\textquotedblleft}
bis zur Krankenanstalt bestimmen, in der Krankenanstalt
wird von 2 Fachärzten der Psychiatrie die eigentliche
{\quotedblbase}Einweisung{\textquotedblleft} verf\"ugt. Während
der Untersuchung und dem Transport darf Zwang angewendet werden,
aber nur durch die gelindesten zur Erfolgsherbeif\"uhrung m\"oglichen
Mittel bei Beachtung der Mittel-Ziel-Relation (=\,Verhältnismä{\ss}igkeitsgrundsatz).

\end{itemize}
\paragraph{Rechtsfolge}
\gFazit{
Alex, Sepp und Bernd haben somit das Recht Gewaltausbr\"uche des Hr.
Hauzu gegen sich und andere abzuwehren. Dabei k\"onnen sie das
verhältnismä{\ss}ig gelindeste Zwangsmittel wählen, sind aber
verpflichtet unverz\"uglich die Polizei und \"uber diese den Amtsarzt
zu verständigen. Trotzdem sollte Gewalt durch Sanitäter immer als
letzte Option verwendet werden, auch wegen dem Risiko sich durch das
Anhalten selbst zu verletzen. Daher der Grundsatz: Selbstschutz geht
vor Fremdschutz (bzw. \"ubertriebenen Heldenmut).
}
\subsubsection{Fall 2. Gefahr durch Hund}
\paragraph{Sachverhalt}

Kaum hat Hr. Hauzu geh\"ort, dass die Polizei und der Amtsarzt
verständigt wurden, st\"urmt er zu einem kleinen Nebenzimmer in
seiner Gemeindebauwohnung, \"offnet die T\"ur und lässt einen
{\quotedblbase}kleinen{\textquotedblleft} (65~kg schweren)
Dobermann-Kampfhund heraus. Mit fletschenden Zähnen nähert er sich
den Sanitätern und droht jeden Moment anzugreifen.

\paragraph{Rechtsfrage}
Was darf gegen den Hund getan werden?

\paragraph{Rechtslage}

\begin{center}
\begin{tabularx}{\textwidth}[H]{XX}
\toprule
\gEs{Rechtfertigender Notstand}
&
\gEs{Entschuldigender Notstand}
\\\midrule
\begin{itemize}
    \item in keinem Gesetz aufgezeichnet; wird aber einheitlich
    angenommen
    \item anzuwenden wenn die Gefahr auch tatsächlich
    (objektiv) vorliegt
\end{itemize}
&
\begin{itemize}
    \item {\S}\,10 StGB; {\S}\,1306a ABGB
    \item anzuwenden wenn man die Gefahr blo{\ss} subjektiv
    annimmt
\end{itemize}
\\\midrule

\gEs{Notstandssituation}: Gefahr ist:

\begin{itemize}
    \item {\mdseries kein menschlicher Angriff}
    \item {\mdseries gegenwärtig oder unmittelbar drohend}
    \item {\mdseries auf beliebiges Rechtsgut gerichtet}
    \item {\mdseries droht bedeutender Nachteil}
\end{itemize}

\gEs{Notstandshandlung}:

\begin{itemize}
    \item {\mdseries Abwehr als schonendstes Mittel}
    \item {\mdseries Zurechnungsprinzip (=~Abwehr geht in Sphäre des
    Gefahrenverursachers)}
    \item {\mdseries geringes Risiko bei hoher Rettungschance }
\end{itemize}
&
\gEs{Notstandssituation}: Gefahr ist:

\begin{itemize}
    \item {\mdseries kein menschlicher Angriff}
    \item {\mdseries gegenwärtig oder unmittelbar drohend}
    \item {\mdseries auf beliebiges Rechtsgut gerichtet}
    \item {\mdseries droht bedeutender Nachteil}
\end{itemize}

\gEs{Notstandshandlung}:

\begin{itemize}
    \item {\mdseries Abwehr als schonendstes Mittel}
    \item {\mdseries Zurechnungsprinzip (= Abwehr geht in Sphäre des
    Gefahrenverursachers)}
    \item {\mdseries Drohender Nachteil darf geringf\"ugig kleiner als das
    Ergebnis der Abwehrhandlung sein.}
\end{itemize}
\\\bottomrule
\end{tabularx}
\end{center}

\paragraph{Rechtsfolge}

\gFazit{
Alex, Sepp und Bernd d\"urfen sich mit den schonendsten Mitteln
verteidigen, um ihr eigenes Leben oder ihre Gesundheit zu sch\"utzen.}

\begin{quotation}
Praxistipp: Schon vor Betreten der Wohnung sollte ein Hund von seinem
Besitzer in einen  separaten Raum weggesperrt werden.
\end{quotation}

\subsection{ Heilbehandlung }

\subsubsection{Fall 3. Blutzuckermessung =  Heilbehandlung }
\paragraph*{Sachverhalt}

Durch die rechtzeitig eintreffende Polizei, kann die Gefahr durch den
Hund beseitigt und Hr. Hauzu festgehaltenen werden. Trotz der Zuweisung
durch den etwas später eingetroffenen Amtsarzt, m\"ochte Bernd,
gewissenhaft wie er ist, einen vollen Neuro-Check mit einem
Blutzuckertest durchf\"uhren.

\paragraph*{Rechtsfrage}
Wann liegt eine  Heilbehandlung  (= HB) vor? Darf der Blutzuckertest
durchgef\"uhrt werden?

\paragraph*{Rechtslage}

\gT{ Heilbehandlung }:
Wird gewonnen aus allgemeinen Rechtsgrundsätzen und der systematischen
Betrachtung der gesamten Rechtsordnung.

\subparagraph*{Voraussetzung:}

\begin{itemize}
 \item medizinisch indiziert (= Heilbehandlung ist zur Diagnose oder Linderung von Erkrankungen und/oder Verletzungen notwendig)
 \item lege artis durchgef\"uhrt (= Heilbehandlung entspricht dem geltenden Stand der Wissenschaft)
 \item durchzuf\"uhren nach Aufklärung des Patienten
 (${\rightarrow}$ siehe ad 4.)
\end{itemize}

${\rightarrow}$ das Vorliegen einer Heilbehandlung wird als Rechtfertigung f\"ur
alle Einwirkungen am Patienten gewertet

${\rightarrow}$ wer die Voraussetzungen der Heilbehandlung nicht einhält macht
sich strafbar nach:

\begin{itemize}
\item Strafrecht:


 \begin{itemize}
  \item a) \gEs{K\"orperverletzung} - {\S}{\S} 83 ff StGB: Sonderfall der HB-Voraussetzungen: es ben\"otigt nur med. Indikation und lege artis Durchf\"uhrung; die Aufklärung ist hierf\"ur irrelevant

  \item b) \gEs{Eigenmächtige  Heilbehandlung } - {\S} 110 StGB: siehe ad 4.

 \end{itemize}

 \item Zivilrecht:

 \begin{itemize}
  \item a) \gEs{Schadenersatzanspr\"uche} - {\S}{\S} 1293 ff ABGB: die Heilbehandlung muss alle drei Voraussetzungen erf\"ullen (auch Aufklärung); Umfang des Schadenersatzes siehe ad 11.
 \end{itemize}

\end{itemize}

\gEs{Blutzuckermessung}:

{\mdseries
Die Messung des Blutzucker mit einer Lanzette war fr\"uher umstritten,
da manche Rettungsorganisationen davon ausgingen, dass es sich um einen
nicht dem Sanitäter erlaubten Eingriff am menschlichen K\"orper
handelt. Dieser Streit wurde durch eine Novellierung des SanG gel\"ost,
in dem man die Blutzuckermessung f\"ur den Sanitäter ausdr\"ucklich
zu lies.}

\paragraph*{Rechtsfolge}
\gFazit{
	Da ein kompletter Neuro-Check bei einem neurologisch
	auffälligen Patienten medizinisch indiziert und
	gesetzlich erlaubt ist, darf jeder der drei Sanitäter eine solche Ma{\ss}nahme lege artis durchf\"uhren. Eine
	Aufklärung ist aufgrund der mangelnden Einsichts- und
	Urteilsfähigkeit nicht m\"oglich.}

\subsection{Selbstbestimmungsrecht}

\subsubsection{Fall 4. Aufklärung und Einwilligung}

\paragraph*{Sachverhalt}

Nachdem nun Hr. Hauzu gebändigt, zur Gänze untersucht und
eingewiesen wurde, freut sich die RTW-Mannschaft schon auf den
wohlverdienten Kaffee am St\"utzpunkt. Doch wie es kommen musste,
k\"onnen Alex, Sepp und Bernd nicht einr\"ucken, da sie einen
Folgetransport zugewiesen bekommen. Die vorgefundene
Patientin Fr. Uninformiert ist eine in der 30. Wo.
schwangere Mutter, die erst letzte Woche ihren 18. Geburtstag feierte. Grund f\"ur ihren Anruf sind starke
Kopfschmerzen und leichte Seh- und H\"oreinschränkungen. Bei der
weiteren Untersuchung durch Sepp werden geschwollene Kn\"ochel, ein
Blutdruck von 140/90 mm/Hg und nach beiläufiger Aussage der
Patientin starker Harndrang mit s\"u{\ss}lichem Geruch
diagnostiziert.

Schnell ist Sepp klar, dass es sich um eine \index{EPH-Gestose} EPH-Gestose handelt, die zur
Eklampsie f\"uhren kann und einen raschen, aber m\"oglichst reizarmen
Transport auf die Gynäkologie ben\"otigt. Da er aber wenig von jungen
M\"uttern hält, erachtet er es auch nicht f\"ur notwendig Fr.
Uninformiert \"uber die nächsten Geschehnisse zu informieren, sondern
ordnet nur harsch Befehle an, die von Fr. Uninformiert
{\quotedblbase}gefälligst{\textquotedblleft} durchzuf\"uhren sind. Er
redet \"uberhaupt bis zur Aufnahme nur das n\"otigste mit ihr.

\paragraph*{Rechtsfrage}
Besteht eine Aufklärungspflicht und was umfasst sie? Muss eine
Einwilligung erfolgen und wenn ja, in welcher Form?

\paragraph*{Rechtslage}

%\begin{multicols}{2}
\subparagraph*{\gEs{Aufklärungspflicht}:}

\begin{itemize}
	\item {\mdseries
	Jeder Patient ist aufzuklären,}

	\item {\mdseries
	vor Setzung der Ma{\ss}nahme,}

	\item wenn er einsichts- und urteilsfähig ist.

\end{itemize}

\subparagraph*{\gEs{Einsichts- und Urteilsfähigkeit}:}

	\begin{itemize}
	\item {\mdseries
	grundsätzlich: bei jedem volljährigen Patienten gegeben }

	\item {\mdseries
	\gE{fehlend bei}: }


	\begin{enumerate}
		\item grundsätzlich unm\"undigen Minderjährigen (0-14 LJ.), wobei es vorrangig auf die Einsichts- und Urteilsfähigkeit ankommt

		\item bei m\"undig Minderjährigen (14.-18 LJ.) vorhanden, solange entsprechende Einsichts- und Urteilsfähigkeit nicht fehlt

		\item psychisch erkrankter Patient mit fehlender Einsichts- und Urteilsfähigkeit

		\item Patient im Schockzustand \gAuthNotes{GABS: CAVE
		Verwechslung mit "`echtem"' Schock -- anders Formulieren?}

		\item stark berauschter Patient (ab ca. 2,5{\textperthousand})

		\item Bewusstlosen

	\end{enumerate}
\end{itemize}
\begin{itemize}
	\item[] ad 1) \& 2): {\S} 146c ABGB:


	\begin{itemize}
		\item grundsätzlich entscheidend ist die vom Einsatzpersonal getroffene
		Beurteilung der Einsichts- und Urteilsfähigkeit des Minderjährigen

		\item nur im Zweifel ist erst der m\"undige Minderjährige selbstständig
		entscheidungsbefugt (es ben\"otigt keine Zustimmung der Eltern) 

		\item nur bei schwerwiegenden Angelegenheiten muss ein Elternteil zustimmen
		${\rightarrow}$ stimmt ein Elternteil zu, der andere verweigert aber,
		so ist  entweder jener Teil entscheidungsbefugt der die gr\"o{\ss}ere
		Rolle in der Erziehung des Patienten spielt (im Zweifel
		die Mutter) oder das \index{Pflegschaftsgericht}Pflegschaftsgericht ist anzurufen

		\item lehnen beide ab, obwohl es sich um eine notwendige Heilbehandlung handelt, muss das
		Pflegschaftsgericht angerufen werden
\end{itemize}
\end{itemize}
\subparagraph*{\gEs{Umfang der Aufklärung}:}

\begin{itemize}
	\item 	Grundsatz: Die Aufklärung hat in jenem Rahmen zu erfolgen, der einer
	ma{\ss}gerechten Figur unter Ber\"ucksichtigung der individuellen
	Besonderheiten des Patienten und der Situation gerecht
	wird.

	\item Detail:
	

	\begin{itemize}
	
		\item Umfang und Tragweite der  Heilbehandlung 
		\item Folgen, Gefahren und Nebenwirkungen der
		 Heilbehandlung 
		\item Umstände die f\"ur verständigen Patienten ins
		Gewicht fallen w\"urden

		\item umgekehrte Proportion zur Dringlichkeit der
		 Heilbehandlung 

	\end{itemize}
\end{itemize}

\subparagraph*{\gEs{Unterlassen der Aufklärung}:}
Unterbleibt die Aufklärung, wird angenommen, dass dem
Patient die Tragweite der Ma{\ss}nahmen nicht bewusst ist und er daher nicht in die
Heilbehandlung einwilligt. 

\subparagraph*{\gEs{Besonderheit }\gT{\index{Revers}Revers}:}

\begin{itemize}
	\item Was ist der \gT{Revers}?
	\begin{itemize}
	\item {\mdseries
	Der Revers ist eine schriftliche Dokumentierung \"uber die Aussage des
	Patient, eine vom Einsatz\-personal empfohlenen Ma{\ss}nahme
	abzulehnen. Die Verweigerung kann auch vom zuständigen Vertreter des
	Patienten abgegeben werden.}
	\end{itemize}	
\end{itemize}
\begin{itemize}
	\item Welche Voraussetzungen hat der Revers?
	\begin{itemize}
		\item auszuf\"ullen nach umfassender Aufklärung des Patienten
		\item Schriftlichkeit (Transportschein, Reversschein oder beigelegtes
		Dokument)
		\item Unterschrift des Patienten (oder seines zuständigen Vertreters)
	\end{itemize}
\end{itemize}
\begin{itemize}
	\item Wann und warum ist der Revers notwendig?
	\begin{itemize}
		\item {\mdseries
		 wenn das Einsatzpersonal eine Ma{\ss}nahme f\"ur notwendig erachtet,
		der Patienten  (oder sein zuständiger Vertreter) diese Ma{\ss}nahme
		aber verweigert}
		\item zur Einhaltung der gesetzlichen Dokumentationspflicht und zur Abwendung
		von  klagsweise geltendmachbaren Anspr\"uchen des Patienten
	\end{itemize}
\end{itemize}
%\end{multicols}


\paragraph*{Rechtsfolgen}

\gFazit{
Aufgrund der unterlassenen Aufklärung machte sich Sepp strafbar nach
folgenden Bestimmungen:

\begin{itemize}
	\item \gEs{Strafrecht}:
	\begin{itemize}
		\item Eigenmächtige  Heilbehandlung  ({\S} 110 StGB): 
		\begin{itemize}
			\item nur bei {\quotedblbase}Gefahr in Verzug{\textquotedblleft}\marginlabel{Gefahr in Verzug: \gE{unmittelbar} drohende oder
			gegenwärtige Gefahr, die ein Abwarten nicht zulässt bzw.
			unmittelbares Handeln verlangt} ben\"otigt es
			keine Aufklärung bzw. Einwilligung und es d\"urfen alle nach Ansicht
			der Einsatzkräfte direkt notwendigen Ma{\ss}nahmen gesetzt werden
		\end{itemize}
		\item K\"orperverletzung \gZ{({\S}{\S} 83 ff StGB)}:
		\begin{itemize}
			\item keine Strafbarkeit, solange die Heilbehandlung medizinische indiziert und lege artis
			durchgef\"uhrt wurde
		\end{itemize}
	\end{itemize}
	\item \gEs{Zivilrecht}:
	\begin{itemize}
		\item 
		Schadenersatz wegen K\"orperverletzung \gZ{({\S}{\S} 1293 ff ABGB)}:

		\begin{itemize}
			\item {f\"ur das Zivilrecht ist die Einwilligung aufgrund einer Aufklärung
			notwendige Voraussetzung f\"ur die rechtfertigende Heilbehandlung ${\rightarrow}$
			mangelt es an ihr, kann Schadenersatz verlangt werden}

		\end{itemize}
		
	\end{itemize}
\end{itemize}
}

\subsubsection{Fall 5.
Patientenverf\"ugung}\label{topic:patientenverfuegung}
\paragraph*{Sachverhalt}

Nach Ablieferung der Fr. Uninformiert drängt Alex seine Kollegen zu
einer schnellen Abfahrt, um endlich seinen Kaffee zu bekommen. Doch
macht ihm der Journaler schon wieder einen Strich durch die Rechnung,
als er das Team zu einem weiteren BO beruft.

Dieser liegt in einem kleinen Pflegeheim, indem eine aufgeregte
Heimpflegerin auf Alex, Sepp und Bernd zust\"urmt und sie zu der 70
jährigen, regungslosen Patientin, Fr. Finito, bringt.
Gleichzeitig hält die Heimpflegerin den Sanitätern eine Patientenverf\"ugung vor, die
Fr. Finito vor \"uber 5 Jahren auf einer Serviette unterschrieben hat
und nach der {\quotedblbase}bei ausbleibenden Kreislauf keine
Wiederbelebungsma{\ss}nahmen gesetzt werden{\textquotedblleft}
d\"urfen. Deshalb hatte die Pflegerin es bis jetzt auch unterlassen mit
einer Reanimation zu beginnen. Da sie aber ebenso nicht wei{\ss}, ob
die Verf\"ugung bindend ist, rief sie die Rettung und setzte sich
wartend neben die Patientin.

\paragraph*{Rechtsfrage}

Was ist eine Patientenverf\"ugung? Ab wann und wie weit ist sie bindend?

\paragraph*{Rechtslage}

%\begin{multicols}{2}
\subparagraph*{\gEs{Allgemeines zur Patientenverf\"ugung}:}

\begin{itemize}
	\item 
	Willenserklärung einer einsicht-, urteils- und
	äu{\ss}erungsfähigen Person

	\item {\mdseries
	erst von Bedeutung, wenn der Patient seinen Willen nicht mehr ausdr\"ucken
	kann}

	\item {\mdseries
	nur Ablehnung einer bestimmten Heilbehandlung kann Inhalt der Verf\"ugung sein}

	\item unbeachtlicher Inhalt:

	\begin{itemize}
		\item Ausschluss einer Grundversorgung

		\item med. nicht indizierte Behandlung
	\end{itemize}
\end{itemize}

\subparagraph*{\gEs{zwei Arten der Patientenverf\"ugung}:}

\begin{itemize}
	\item \gT{Verbindliche Patientenverf\"ugung}:


	\begin{itemize}
		\item schriftlich vor Notar, Rechtsanwalt oder rechtskundigen
		Patientenvertreter

		\item eigenhändig unterschrieben und datiert

		\item aufgrund einer umfassenden Aufklärungen durch einen Arzt

		\item Inhalt muss eine konkrete Heilbehandlung ablehnen ${\rightarrow}$ allgemeine oder
		widerspr\"uchliche Formulierungen f\"uhren zur Unwirksamkeit

		\item Wirkung nur f\"ur 5 Jahre

		\item Eine Verf\"ugung die alle Elemente enthält ist absolut bindend und
		darf nicht umgangen werden
	\end{itemize}

	\item \gT{Beachtliche Patientenverf\"ugung}:

	\begin{itemize}
		\item {\mdseries
		\gEs{grundsätzlich}: alle Verf\"ugungen, die nicht
		die obigen Voraussetzungen erf\"ullen}

		\item {\mdseries
		\gEs{Umfang der Beachtlichkeit}: blo{\ss} relative
		Beachtlichkeit:}

		\begin{itemize}
			\item {\mdseries
			Verf\"ugung soll in Behandlungsentscheidungen ber\"ucksichtigt werden,
			aber}

			\item {\mdseries
			das Wohl des Patienten bleibt oberstes Gebot }

			\item {\mdseries
			daher kann man sich bei besonderer Begr\"undung \"uber die Verf\"ugung
			hinwegsetzen}

		\end{itemize}
	\end{itemize}
\end{itemize}

\begin{itemize}
	\item {\mdseries
	\gEs{Absolute Unwirksamkeit/Unbeachtlichkeit bei}:}


	\begin{itemize}
		\item {\mdseries
		zu allgemeine oder widerspr\"uchliche Formulierungen}

		\item {\mdseries
		rechtswidriger Inhalt}

		\item {\mdseries
		erkennbare Willensmängel des Patienten bei Abgabe der
		Verf\"ugung (nicht ernst gemeint, Irrtum, List, Drohung oder Täuschung)}
	\end{itemize}
\end{itemize}

\subparagraph*{\gEs{Suche nach Patientenverf\"ugung} --
\gZ{{\S} 12 PatVG:}}~

Bei Notfallversorgungen muss nicht nach einer Verf\"ugung gesucht
werden, wenn die Suche das Leben oder die Gesundheit des Pat ernstlich
gefährden w\"urde. Daher wird f\"ur das Rettungswesen idR. keine
Verpflichtung zur Suche, sondern zur Behandlung bestehen.

%\end{multicols}

\paragraph*{Rechtsfolge}

\gFazit{
	Die Patientenverf\"ugung gilt mit \"uber 5 Jahren, der fehlenden
	ärztlichen Aufklärung und dem mangelnden Beiziehen eines
	Rechtsverständigen nicht mehr als verbindlich und wäre somit als
	blo{\ss} beachtlich einzustufen. 


	Doch aufgrund des zu allgemeinen Inhalts, der Beiläufigkeit der
	Verf\"ugung (Serviette als Unterlage) und wegen des hohen Alters der
	Verf\"ugung kann kein konkreter und aktueller Wille der
	Patientin erkannt werden, was dazu f\"uhrt, dass das Wohl
	der Patientin das oberste Gebot bleibt. 
	
	Mit einer Reanimation ist zu beginnen.

}

\subsection{Rechte und Pflichten des Sanitäters}

\subsubsection{Fall 6. Unterlassung der Notarztnachforderung}

\paragraph*{Sachverhalt}

Endlich am St\"utzpunkt angekommen, ist Alex auch schon der Erste der
sich zum Kaffeeautomaten begibt. Entkräftet durch die einsetzende
{\quotedblbase}Hypokoffeinomie{\textquotedblleft} kramt er nach
Kleingeld und kaum als er es einwurfbereit in Händen hält, t\"ont
die nächste Ausfahrt durch die Lautsprecher, mit dem Hinweis einer
dringenden Berufung. Alex geht aber in seiner typisch ruhigen Manie,
gl\"ucklich durch die enteral erlangte Sedierung, zum Fahrzeug, mit dem
Spruch auf den Lippen: {\quotedblbase}immerhin konn des
Wagl{\textquoteright} eh net ohne mie foahrn!{\textquotedblleft}

Nach nur kurzer Anfahrtszeit begegnet Alex, Sepp und Bernd im ersten
Stock eines kleinen Reihenhaus der 55 jährige, adip\"ose Hr. Schmerz
mit Schmerzen hinter der Brust, die in den linken Arm ausstrahlen und
ihm selbst in Ruhelage mit Atemnot und kalten Schwei{\ss}ausbr\"uchen
plagen. F\"ur Sepp ist sofort klar, dass es sich um einen
MCI\footnote{\gE{MCI}: Myocardinfarkt: Herzinfarkt.} handelt.
Er m\"ochte schleunigst einen NEF nachfordern, doch Alex
meint, dass er keine Lust hätte schon wieder auf
einen {\quotedblbase}Boder\footnote{"Bader": fr\"uher
Kundiger der Heilbäder, heute umgangssprl. f\"ur
Arzt}{\textquotedblleft} zu warten und der Patient auch
selbst runter zum Fahrzeug gehen sollte, da er es im Kreuz
hätte und sowieso noch eine ca. 20 Minuten lange
Anfahrtszeit zur Krankenanstalt durch den schon einsetzende Morgenstau
bewältigen m\"usste.

\paragraph*{Rechtsfrage}

Besteht die Pflicht f\"ur Sanitäter zur Nachforderung eines Notarztes?

\paragraph*{Rechtslage}

\gEs{Unterlassene Hilfeleistung} - {\S} 95 StGB

\begin{itemize}
	\item {\mdseries
	jeder Durchschnittsmensch ist zur Durchf\"uhrung von
	\gE{ihm zumutbaren Hilfeleistungen} verpflichtet. Die
	Zumutbarkeit bestimmt sich dabei nach dem individuellen Stand an
	Ausbildung, Fähigkeiten und Kenntnissen. Aber schon zu den Grundlagen
	einer Ersten Hilfe, die von jedem Menschen mit Besuch eines
	F\"uhrerscheinkurses erwartet werden kann, geh\"ort das Absetzen eines
	Notrufes zur  Anforderung von geeigneten Rettungskräften.}

\end{itemize}

\gEs{{\S} 4 SanG}

\begin{itemize}

	\item {\mdseries
	ein jeder Sanitäter, unabhängig von seinem Ausbildungsstand, ist zur
	Durchf\"uhrung von lebensrettenden Sofortma{\ss}nahmen verpflichtet,
	was gleichzeitig beinhaltet, dass er bei Notwendigkeit eine
	unverz\"ugliche Anforderung des Notarztes oder sonstigen Arztes zu
	veranlassen hat.}

\end{itemize}

\paragraph*{Rechtsfolge}

\gFazit{
	W\"urde Alex die zweckmä{\ss}ige Nachforderung eines Notarztes unterlassen,
	macht er sich gleichzeitig nach strafrechtlichen ({\S} 95 StGB) und
	verwaltungsrechtlichen Bestimmungen ({\S} 9 SanG) strafbar.
}

\subsubsection{Fall 7. Sanitäterkompetenzen}

\paragraph*{Sachverhalt}

Nach langem hin und her, k\"onnen Sepp und Bernd gemeinsam den Alex
\"uberreden doch einen Notarzt nachzufordern und Sepp, der gerade frisch
gebackener NFS geworden ist, sieht endlich seine M\"oglichkeit
umfangreich am Hr. Schmerz zu werken. Dabei m\"ochte er neben der
Verabreichung von Nitrolingual-Spray auch gleich einen ven\"osen Zugang
legen, da er das sowieso als ehemaliger Fleischhauer besser als jeder
Arzt kann.

\paragraph*{Rechtsfrage}

Was sind die \gT{Kompetenzen} der einzelnen Sanitäterstufen
(RS/NFS/NKA/NKV/NKI)? Darf man die Kompetenzgrenzen (bei \gE{Gefahr in
Verzug}) \"uberschreiten?

\paragraph*{Rechtslage}

\gEs{Kompetenzstufen} gem. \gZ{{\S}{\S} 9-12} SanG:

\begin{itemize}
	\item {\mdseries
	\gT{RS} - \gZ{{\S} 9} SanG:}

	\begin{itemize}
		\item {\mdseries
		Z 1: selbstständige und eigenverantwortliche Versorgung und Betreuung
		des Patienten}

		\item {\mdseries
		Z 2: \"Ubernahme und \"Ubergabe des Patienten im Rahmen
		eines Transportes}

		\item {\mdseries
		Z 3: Hilfestellung bei Akutsituationen; Verabreichung von O2}

		\item {\mdseries
		Z 4: Durchf\"uhrung von lebensrettenden Sofortma{\ss}nahmen}

		\item {\mdseries
		Z 5: Durchf\"uhrung von Sondertransporten}

		\item {\mdseries
		Z 6: Blutzuckermessung}

	\end{itemize}

	\item \gT{NFS} - \gZ{{\S} 10} SanG:

	\begin{itemize}
		\item Z 1: Tätigkeiten gem. {\S} 9
		\item Z 2: Unterst\"utzung des Arztes; Betreuung von Transporten
		\item 
		Z 3: Verabreichung der \gEs{Arzneimittelliste 1}: 
		Wird vom leitenden Arzt der Rettungsorganisation
		festgelegt . Die Verabreichung kann
		durchgef\"uhrt werden ohne dass ein Notarzt informiert werden muss oder es
		ein solcher anordnet.
		\item Z 4: eigenverantwortliche Betreuung von berufsspezifischen
		Medizinprodukten und Arzneimitteln
		\item 	Z 5: Mitarbeit in Forschung
	\end{itemize}
	\item \gT{NKA} - \gZ{{\S} 11 Abs. 1 Z 1} SanG:
	zusätzlich zu den oben genannten Kompetenzen darf der NKA jene
	Medikamente der Arzneimittelliste 2 verabreichen, die ohne ven\"osen
	Zugang verabreicht werden k\"onnen.
	
	Die \gEs{Arzneimittelliste 2} wird vom leitenden
		Arzt der Rettungsorganisation festgelegt und darf erst
		verabreicht werden:
			\begin{itemize}
		\item auf Anweisung durch anwesenden Arzt (jeder mit jus
		practicandi -- nicht unbedingt Notarzt) oder
		\item
		nach Verständigung eines Notarztes.
	\end{itemize}
	\item 
	\gT{NKV} - \gZ{{\S} 11 Abs. 1 Z 2} SanG:
	zusätzlich zu den oben genannten Kompetenzen darf der NKV jene
	Medikamenten der Arzneimittelliste 2 verabreichen, die \"uber ven\"osen
	Zugang verabreicht werden k\"onnen.
	\item 	\gT{NKI} - \gZ{{\S} 12 SanG}:
	Durchf\"uhrung einer endotrachealen Intubation samt Beatmung, aber ohne
	Prämedikation
	\begin{itemize}
		\item {\mdseries
		NKI-Kompetenz muss alle 2 Jahre rezertifiziert werden}

		\item {\mdseries
		Durchf\"uhrung der Intubation erst:}

		\begin{itemize}
			\item {\mdseries
			auf Anweisung durch anwesenden Arzt (jeder mit jus practicandi -- nicht unbedingt Notarzt) oder}

			\item {\mdseries
			nach Verständigung eines Notarztes }

		\end{itemize}
	\end{itemize}
\end{itemize}	
{\mdseries
\gEs{\"Uberschreiten der Kompetenzstufen}:}

Die Durchf\"uhrung von Ma{\ss}nahmen au{\ss}erhalb des eigenen
Kompetenzstandes, f\"uhrt zu einer Verwaltungsstrafe nach SanG bzw.
einer gerichtlichen Strafe wegen K\"orperverletzung und
zivilrechtlichen SchEA. Selbst bei Wissen \"uber Gefahren, Risiken und
Durchf\"uhrungsart einer kompetenzmä{\ss}ig h\"oheren Ma{\ss}nahme,
gibt es keine Rechtfertigung durch {\S} 95 StGB (Unterlassene
Hilfeleistung).

Dies bedeutet weiters, dass einfache Ma{\ss}nahmen, die zwar das
Berufsgesetz nicht aufzählt, aber von jedem Laien (mitunter nach
kurzer Einweisung) durchgef\"uhrt werden d\"urfen, auch von
Sanitätern gesetzt werden d\"urfen bzw. iSd. {\S} 95 StGB gesetzt
werden m\"ussen.

\paragraph*{Rechtsfolge}

\gFazit{
Sepp darf keinen ven\"osen Zugang legen oder sogar Medikamente dar\"uber
verabreichen, selbst  wenn Gefahr in Verzug vorliegt. }


\subsubsection{Fall 8. Rechtswidrige Anordnung durch Notarzt}

\paragraph*{Sachverhalt}

Sepp konnte sich nun dazu \"uberreden lassen keinen Venflon zu legen und
m\"ochte wenigstens Nitro verabreichen. Gerade noch rechtzeitig bemerkt
er eine offene Packung von Viagra am Tisch liegen. Hr. Schmerz meint,
als er Sepp die Packung betrachten sieht, dass das
{\quotedblbase}Zeug{\textquotedblleft} nichts nutzt, da er vor zwei
Stunden eine Tablette eingenommen hätte und noch immer
{\quotedblbase}nichts steht{\textquotedblleft}. Kurze Zeit später
trifft der NEF mit der Notärztin Natascha ein, die den
Patienten flapsig untersucht und eine sofortige
Nitro-Verabreichung zur Bekämpfung des hohen Blutdrucks (200/110\,mmHg)
anordnet, obwohl sie von Sepp \"uber die gefundene Viagra-Packung informiert wurde.

\paragraph*{Rechtsfrage}

Ist der Sanitäter verpflichtet eine rechtswidrige bzw. gefährliche
Anordnung des Notarztes (oder eines anderen Vorgesetzten) zu folgen?

\paragraph*{Rechtslage}

\gEs{Ablehnung von rechtswidrigen Anordnungen}:

\begin{itemize}
	\item {\mdseries
	Zivildiener: {\S} 22 Abs. 2 Zivildienstgesetz:}


	\begin{itemize}
		\item
		{\quotedblbase}Er darf die Befolgung einer Weisung nur dann ablehnen,
		wenn [{\dots}] die Befolgung gegen strafgesetzliche Vorschriften
		versto{\ss}en w\"urde.{\textquotedblleft}

	\end{itemize}

	\item {\mdseries
	Haupt- und Ehrenamtliche: }

	\begin{itemize}
		\item {\mdseries
		{\S} 4 Abs. 1 S. 2 SanG}

		\item {\mdseries
		{\S} 7 iVm {\S} 879 Abs. 1 ABGB}

		\item {\mdseries
		{\S} 22 Abs. 2 ZDG analog:}

	\end{itemize}

\end{itemize}

Aufgrund der im SanG fehlenden konkreten Bestimmung eines
Ablehnungsrechtes von Sanitätern, kann ein solches nur kompliziert
gewonnen werden aus:

\begin{itemize}
	\item {\mdseries
	allgemeinen Verpflichtung zur Wahrung des Patientenwohls nach {\S} 4
	SanG}

	\item {\mdseries
	nat\"urlichen Rechtsgrundsätzen iSd. {\S} 7 ABGB und dem
	Nichtigkeitsgebot bei Sittenwidrigkeit des {\S} 879 ABGB}

	\item {\mdseries
	analoger Anwendung des Zivildienstgesetz (ZDG) }

\end{itemize}

\paragraph*{Rechtsfolge}

\gFazit{
	Die RTW-Mannschaft darf daher die f\"ur den Patiente
n	schädigende Anordnung ablehnen. Natascha hingegen macht
	sich strafbar wegen K\"orperverletzung und riskiert uU. zivilrechtliche SchEA
	(${\rightarrow}$ Nitro-Gabe nicht lege artis).}


\subsubsection{Fall 9. Dokumentationspflicht}

\paragraph*{Sachverhalt}

Genervt von der Auseinandersetzung der Notärztin mit den Sanitätern,
beginnt Hr. Schmerz wild vor Schmerzen die Beteiligten anzubr\"ullen,
dass sie doch endlich ihre {\quotedblbase}gottverdammte
Arbeit{\textquotedblleft} tun sollen. Der ebenso gereizte Bernd belehrt
daraufhin Hr. Schmerz, dass er sich doch gefälligst wie ein
{\quotedblbase}richtiger Patient{\textquotedblleft} verhalten soll und
sich erst r\"uhren darf, wenn es ihm gesagt wird. Diese Aussage wird
Bernd aber zum Verhängnis, da sich Hr. Schmerz Wochen später beim
Arbeitgeber von Alex, Sepp und Bernd \"uber das
{\quotedblbase}beleidigende Verhalten{\textquotedblleft} des Bernd
beschwert. Sepp seinerseits ist \"uberzeugt, dass er sich korrekt
verhalten hat, doch kann er sich so recht an den alten Fall nicht mehr
erinnern.

\paragraph*{Rechtsfrage}

Besteht eine Dokumentationspflicht und wenn ja, in welchem Umfang? 

\paragraph*{Rechtslage}

\gZ{{\S}\,5} \gEs{SanG}:
{\quotedblbase}\gC{Sanitäter haben [{\dots}] die von
ihnen gesetzten Ma{\ss}nahmen zu dokumentieren.}{\textquotedblleft}

\gEs{Vertragsrechtliche Relevanz}:

Im \"osterreichischen Recht hat grundsätzlich der Kläger den
Schaden, die Rechtswidrigkeit und das Verschulden des Beklagten zu
beweisen. Besteht aber zwischen Kläger und Beklagten ein Vertrag, so
wird gem. \gZ{{\S}\,1298} AGBG das Verschulden des Beklagten
angenommen bzw. der Beklagte muss beweisen, dass ihn kein Verschulden trifft (=
\index{Beweislastumkehr}\gT{Beweislastumkehr}). 

Im Rettungswesen besteht zu jedem Patient ein \index{Behandlungsvertrag}\gT{Behandlungsvertrag}, der eine
solche Beweislastumkehr bewirkt. Um nun f\"ur den Beklagten die
M\"oglichkeit zu schaffen, sein Unverschulden zu beweisen, wurde die
Dokumentationspflicht eingef\"uhrt, mit der bei ausreichender
Dokumentation wieder der Kläger das Verschulden des Beklagten zu
beweisen hat.

{\mdseries
\gEs{Inhalt der Dokumentation}:}

\begin{itemize}
	\item {\mdseries
	gesetzte med. Ma{\ss}nahmen}

	\item {\mdseries
	situationsrelevante Umstände (beim Patient; am BO; während des
	Transportes; {\dots})}

	\item {\mdseries
	Verweigerungen oder besondere Anmerkungen des Patiente
n	(kein Gurt; Marcoumarpflichtig; {\dots})}

	\item {\mdseries
	uU. relevante Inhalte der Aufklärung}

	\item {\mdseries
	sonstige f\"ur den Sanitäter entscheidend wirkende Umstände }

\end{itemize}
\begin{quotation}
Grundsatz: Alles darf dokumentiert werden, es darf aber nicht Nichts
dokumentiert werden. 
\end{quotation}

\subparagraph*{Durchf\"uhrung der Dokumentation:}
Die Dokumentation soll vorrangig am Transportschein selbst oder wenn
nicht mehr ausreichend auf einem beigelegten Blatt Papier erfolgen.
Verweigert der Patient medizinische oder sonst notwendige
Ma{\ss}nahmen, ist von ihm ein Revers zu unterschreiben, der dem
Transportschein beizulegen ist. Dabei hat der Sanitäter
immer (!) durch ein Gespräch auf m\"ogliche Risiken und Gefahren der
Verweigerung hinzuweisen.

Des Weiteren ist eine zusätzliche eigene Dokumentation f\"ur jeden
Sanitäter empfehlenswert.

\paragraph*{Rechtsfolge}

\gFazit{
	Bernd hat keine ausreichende Dokumentation durchgef\"uhrt und kann sich
	somit nicht gegen die Beschuldigungen zur Wehr setzen.}


\subsubsection{Fall 10. Verschwiegenheits- und Auskunftspflicht}

\paragraph*{Sachverhalt}

Durch die Aufregung verschlechtert sich allmählich der Zustand des Hr.
Schmerz und das RTW-Team samt NEF-Mannschaft sind sich einig, dass ein
schleunigster Abtransport durchzuf\"uhren ist. Noch bevor sie aber den
Patienten umlagern k\"onnen, erscheint Fr. Aufgeregt, die
Angeh\"orige des Hr. Schmerz, und verlangt \"uber alles informiert zu werden.

\paragraph*{Rechtsfrage}

Was umfasst die Verschwiegenheitspflicht der Sanitäter und \"Arzte?
Besteht auch eine Auskunftspflicht derselben?

\paragraph*{Rechtslage}

Verschwiegenheitspflicht:

\begin{itemize}
	\item {\mdseries
	\gEs{strafrechtliche} Norm: \gZ{{\S}~121} StGB}

	\item {\mdseries
	\gEs{verwaltungsrechtliche} Norm: \gZ{{\S}~6}~SanG}

	\item {\mdseries
	\gE{beiden gleich}: Geheimnisse, die man aufgrund seiner
	Tätigkeit als Arzt/Sanitäter anvertraut bekommen hat, d\"urfen
	nicht an unbefugte Dritte weitergegeben werden}

	\begin{itemize}
	
		\item {\mdseries
		Geheimnisse: alle Informationen, die die Person des
		Patienten, seinen Gesundheitszustand oder die Umstände am
		Einsatzort betreffen}

		\item {\mdseries
		daher darf man mit Kollegen oder sonstigen Personen \"uber die erlebten
		Geschehnisse reden, solange aus dem Gesagten nicht erkennbar ist, um
		welche Person es sich bei dem Patienten handelt }

		\item {\mdseries
		Bsp: Name des Patienten, Versicherungsnummer, Wohnadresse,
		{\dots}}
		
	\end{itemize}
	
\end{itemize}

\gEs{Wann \gE{darf} man Geheimnisse weitergeben?}

\begin{itemize}
	\item {\mdseries
	Identität des Patienten bleibt gewahrt}

	\item {\mdseries
	bei Entbindung von der Schweigepflicht durch den Patienten}

\end{itemize}

\gEs{Wann \gE{muss} man Geheimnisse weitergeben?} = \gT{Auskunftspflicht} \index{Auskunftspflicht} --
\gZ{{\S} 7} SanG

\begin{itemize}
	\item {\mdseries
	wenn es zur Behandlung des Patienten dient}

	\item {\mdseries
	wenn es gesetzlichen vorgeschrieben ist (Bsp: im Gerichtsverfahren als
	Zeuge)}

	\item {\mdseries
	zum Schutz von h\"oherwertigen Interessen der \"offentlichen
	Gesundheitspflege oder Rechtspflege (${\rightarrow}$ zur Meldepflicht
	siehe ad G)}

\end{itemize}

\gEs{Wem gegen\"uber muss man Auskunft erteilen?}

\begin{itemize}
	\item {\mdseries
	Patienten}

	\item {\mdseries
	gesetzlichen Vertreter des Patienten}

	\item {\mdseries
	anderen Personen der Gesundheitsberufe, die den Patiente
n	betreuen, behandeln oder pflegen bzw. an denen der Patient
	im Zuge des Transportes \"ubergeben wird}

	\item {\mdseries
	Sozialversicherung und Krankenanstalt}

	\item {\mdseries
	vom Patient als auskunftsberechtigt benannte Personen}

\end{itemize}

\paragraph{Rechtsfolge}

\gFazit{
	Fr. Aufgeregt zählt als blo{\ss}e Angeh\"orige nicht zum Kreis der
	Auskunftsberechtigten, solange der Patient die Weitergabe von
	Informationen nicht gestattet. }


\subsubsection{Fall 11. Fortbildungs- und Rezertifizierungspflicht}

\paragraph*{Sachverhalt}

Nachdem nun auch Fr. Aufgeregt beruhigt wurde, wird Hr. Schmerz auf den
Sessel umgelagert, um \"uber die enge Treppe den ersten Stock zu
verlassen. Doch handelt es sich beim Tragsessel um einen neuen
Striker-Sessel, der dem Bernd noch vollkommen neu ist, da es den Sessel
zum Zeitpunkt seiner RS-Ausbildung noch gar nicht gegeben hat. Aber
genau wegen dieser Neuheit hatte sein Arbeitgeber eine interne
Fortbildung \"uber die {\quotedblbase}Neuen
Medizinprodukte{\textquotedblleft} abgehalten, zu der Bernd aber nicht
gekommen ist, da er meinte sowieso schon alles zu wissen und zu keinen
Fortbildungen verpflichtet ist.

\paragraph*{Rechtsfrage}

Besteht eine Fortbildungs- oder sogar Re\-zertifizierungs\-pflicht?

\paragraph*{Rechtslage}

\subparagraph*{\gEs{Fortbildungspflicht nach} \gZ{{\S}~50} \gEs{SanG}:}

Jeder Sanitäter ist dazu verpflichtet sich innerhalb von 2 Jahren ab
seinem Stichtag im Umfang von 16 Stunden fortzubilden. Als Stichtag
wird der Monatserste nach bestandener Rettungssanitäterausbildung
gewertet, wobei sich der Stichtag auch bei Erreichen von h\"oheren
Kompetenzniveaus nicht ändert. Zu welchem Zeitpunkt die 16
Fortbildungsstunden innerhalb des Rahmens von 2 Jahren besucht werden,
ist dem Sanitäter \"uberlassen. 

Man kann also entweder regelmä{\ss}ig Fortbildungen besuchen oder auf
einmal die Stundenzahl ansammeln. Weiters muss ein Sanitäter keine
bestimmten Themen als Fortbildungen besuchen, sondern kann nach seinem
Interesse vorgehen.

\subparagraph*{\gEs{Fortbildungspflicht nach} \gZ{{\S}\,83}
\gEs{MPG}:}

Die im Rettungswesen verwendeten Materialien gelten, abgesehen von
Arzneimitteln, als Medizinprodukte und d\"urfen daher iSd. MPG nur von
eingeschulten Personal angewendet werden. Diese Einschulungen m\"ussen
nicht nur alle Medizinprodukte der Rettungsorganisation beinhalten,
sondern auch von allen Sanitätern besucht werden. Daher kann ein
Sanitäter in Bezug auf Stundenanzahl oder Fortbildungsthemen nicht
frei wählen. Die Fortbildungsstunden f\"ur das MPG k\"onnen aber als
Fortbildungsstunden iSd. SanG gewertet werden. Zusätzlich ist zu
sagen, dass eine Einschulung der meisten Medizinprodukte schon durch
die Standardausbildung zum Rettungs\-sani\-täter abgedeckt ist. F\"ur
die anderen, insbesondere h\"ohere technische Produkte, muss jede
Org\-anisation separate Fortbildungen anbieten.

\subparagraph*{\gEs{Rezertifizierungspflicht} \gEs{gem}.\gZ{{\S}~51} \gEs{SanG}:}

Innerhalb des 2 Jahre Rahmens der Fortbildungspflicht muss sich ein
Sanitäter, unabhängig von seinem Ausbildungsstand, pr\"ufen lassen
in der Herz-Lungen-Wiederbelebung mit halbautomatischen Defibrillator.
Der Rezertifizierungszeitpunkt innerhalb des 2 Jahre Rahmens kann vom
Sanitäter frei gewählt werden.

\gE{Sonderfall NKI}: Ein Notfallsanitäter mit
Kompetenzen der Intubation muss sich zusätzlich innerhalb seines 2
Jahre Rahmens einmal \"uber die Intubation pr\"ufen lassen.

\paragraph*{Rechtsfolge}

\gFazit{
	Die f\"ur den vorliegenden Fall entscheidende Fortbildung gem. MPG wurde
	von Bernd nicht besucht, was ihm nicht nur eine Verwaltungsstrafe,
	sondern auch eine gerichtliche Strafe und zivilrechtlichen SchEA
	einbringen kann. Denn sollte es im Umgang mit dem Medizinprodukt zu
	K\"orperverletzungen kommen, wird aufgrund der mangelnden
	Materialkenntnis die Behandlung als nicht-lege-artis gewertet. Um einer
	solchen Strafe zu entgehen gen\"ugt es, dass Bernd noch vor Ort, aber
	vor Anwendung, durch eine {\quotedblbase}geeignete
	Person{\textquotedblleft} eingeschult wird. 

	Als {\quotedblbase}geeignete Person{\textquotedblleft} gelten vorrangig
	Medizinproduktebeauftragter,  Lehrsanitäter oder Praxisanleiter sowie
	jene, die von der Rettungsorganisation als geeignet genannt werden.

	Da Sepp Praxisanleiter ist, darf er vorab Bernd einweisen.
	}


\subsection{Schadenersatz}

\subsubsection{Fall 12. Verletzung des Patienten durch
Sanitäter:}

\paragraph*{Sachverhalt}

Nachdem nun Sepp unseren Bernd kurzum erklärt hat, wie er mit den
Striker-Sessel umzugehen hat unterlässt er trotzdem das Anschnallen
des Patienten, da er meint, dass {\quotedblbase}eh kana däs
mocht{\textquotedblleft}. Leider ist Hr. Schmerz kein Fliegengewicht
und der Treppenaufgang ist derart verwinkelt, dass es alle M\"uhe und
Not bedarf den Tragsessel ruhig zu halten. 

Fast am Ende des Abgangs angekommen, flitzt auf einmal die
aufgeschreckte Katze des Hr. Schmerz an den beiden
{\quotedblbase}Trägern{\textquotedblleft} vorbei, wodurch der
vorangehende Bernd seine Konzentration verliert, ausrutscht und samt
dem Patient seitlich nach vorne auf einen Kasten und die
darunter sich versteckende Katze namens \emph{"`Muschi"'}
fällt. Bernd kann sich beim Sturz wenigsten selber abst\"utzen und bleibt (auch durch die abfedernde Wirkung der Katze)
unverletzt. Doch Hr. Schmerz prallt mit Schädel und Brustkorb
ungesch\"utzt gegen den Kasten und wird unter dem hinter ihm
nachkommenden Tragsessel begraben. Er zieht sich neben einer RQW auf
der Stirn auch noch mehrere gebrochene Rippen zu. 

\paragraph*{Rechtsfrage - Schadenersatz wegen K\"orperverletzung}

Wann wird man schadenersatzpflichtig? Was umfasst der Schadenersatz? Wer
wird zur Haftung gezogen?

\paragraph*{Rechtslage}

\gEs{Rechtliche Grundlage eines
Schadenersatzanspruches} (= SchEA):

\begin{itemize}

	\item {\mdseries
	allgemein: \gZ{{\S}{\S} 1293 ff ABGB}}

	\item {\mdseries
	konkret f\"ur K\"orperverletzungen: \gZ{{\S}{\S}
	1325-1327 ABGB}}

\end{itemize}

\gEs{Voraussetzungen der Schadenersatzpflicht}:

\begin{itemize}

	\item[] {\mdseries
	1. Schaden = Nachteil am K\"orper, an der Gesundheit, im Verm\"ogen}

	{\mdseries
	2. Kausalität = bei Wegdenken der Schädigerhandlung, bleibt auch der
	Erfolg aus}

	{\mdseries
	3. Rechtswidrigkeit = objektiver Versto{\ss} gegen Gesetze oder
	Vertragspflichten}

	{\mdseries
	4. Schuld = subjektiver Versto{\ss} gegen jene Sorgfaltspflichten, die
	man einhalten hätte k\"onnen}

\end{itemize}

\gEs{Umfang des Schadenersatzes bei
K\"orperverletzungen}:

\begin{itemize}

	\item {\mdseries
	Heilungskosten}

	\item {\mdseries
	Verdienstentgang}

	\item {\mdseries
	Schmerzengeld}

\end{itemize}


\gEs{Wer haftet wie}:

Schädigt der Sanitäter im Zuge seines Dienstes einen Patient oder
dessen Angeh\"origen, kann der Geschädigte nach
\gZ{{\S} 1313a} ABGB direkt gegen die
Rettungsorganisation seine SchEA geltend machen. Denn ein Sanitäter
gilt in dieser Beziehung als Erf\"ullungsgehilfe der
Rettungsorganisation, wodurch er als verlängerter Arm der
Organisation deren vertragliche Pflichten (sorgfältige Behandlung)
erf\"ullt. Daher kann der Geschädigte direkt gegen die Organisation
alle Anspr\"uche wie gegen den Sanitäter selber durchsetzen.

\paragraph*{Rechtsfolge}

\gFazit{
	Im vorliegenden Fall ist der Schaden durch RQW und Frakturen neben der
	Kausalität gegeben. Weiters verst\"o{\ss}t es gegen die gesetzlichen
	und vertraglichen Pflichten den Patient ohne anzuschnallen zu
	transportieren, weshalb eine Rechtswidrigkeit bejaht werden kann. 
	
	Auch hätten sich unsere Sanitäter anders, nämlich den Patient
	anzuschnallen, verhalten k\"onnen, wodurch ebenso die Schuld gegeben
	ist. Daher kann Hr. Schmerz SchEA gegen die Sanitäter oder,
	wirtschaftlich vern\"unftiger, gegen die Rettungsorganisation geltend
	machen.
	}


\subsubsection{Fall 13. Sachbeschädigung}

\paragraph*{Sachverhalt}

Leider sind die Verletzungen am Patient nicht die einzigen Schäden die
entstanden sind, denn der umgeschmissene Kasten war eine
Biedermeier-Rarität mit einem Wert von stolzen 10.000 Euro. Auch die
Katze verwirkte als Auffangkissen f\"ur Bernd ihr Leben und hatte neben
dem Sachwert von 350 Euro auch einen besonderen emotionalen Wert f\"ur
Hr. Schmerz.

\paragraph*{Rechtsfrage: Schadenersatz wegen Sachbeschädigung}

Wann wird man schadenersatzpflichtig? Was umfasst der Schadenersatz? Wer
wird zur Haftung gezogen?

\paragraph*{Rechtslage}

\gEs{Rechtliche Grundlage eines
Schadenersatzanspruches}:

\begin{itemize}
	\item {\mdseries
	allgemein: \gZ{{\S}{\S} 1293 ff ABGB}}

	\item {\mdseries
	konkret f\"ur Sachschäden: \gZ{{\S} 1331 ABGB}}

\end{itemize}

\gEs{Voraussetzungen der Schadenersatzpflicht}:

\begin{itemize}
	\item[] {\mdseries
	gleiche Voraussetzungen wie bei K\"orperverletzungen}

\end{itemize}

\gEs{Umfang des Schadenersatzes bei
Sachbeschädigung}:

\begin{itemize}
	\item {\mdseries
	realer Schaden (= tatsächlich eingetretener Schaden)}

	\item {\mdseries
	entgangener Gewinn (= in Zukunft drohender Schaden)}

\end{itemize}

\paragraph*{Rechtsfolge}

\gFazit{
	gleiche L\"osung wie bei Fall 12.}


\subsection{Behandlungspflicht}

\subsubsection{Fall 14. Behandlungspflicht von Krankenanstalten, \"Arzten,
Sanitätern und Durchschnittsb\"urgern:}

\paragraph*{Sachverhalt}

Vor Schmerzen und Trauer, um seine geliebte Katze
{\quotedblbase}Muschi{\textquotedblleft}, beschimpft und verflucht Hr.
Schmerz die Anwesenden mit wilden Gestikulierungen. Als er dann
wenigstens von Alex und Sepp vorbildlich versorgt wird, beruhigt sich
Hr. Schmerz und lässt sich sogar m\"uhelos ins Auto transportieren. 

Nach ausreichender Versorgung des Hr. Schmerz begibt sich der
RTW auf die Reise zur Krankenanstalt. Doch ist leider das
einzig abbuchbare Bett, welches neben einer Internen- auch
eine Unfall-Ambulanz aufzuweisen hat, im Wilhelminenspital,
welches durch den gro{\ss}en Morgenstau gute 15 Min. entfernt
ist. Schon nach kurzer Fahrzeit klagt Hr. Schmerz \"uber
stärker werdende Atemprobleme und ein rasselndes, brodelndes
Atemgeräusch t\"ont durch den gesamten Patientenraum. Es hat
sich ein kardiales Lungen\"odem gebildet und Natascha f\"uhlt
sich, trotz Verabreichung von Diuretika, nicht mehr in der
Lage den Patient während der Fahrt ausreichend zu behandeln.
Alex soll so schnell wie m\"oglich das nächste Spital
anfahren. Durch eine kurze Abweichung der Route, gelangt das
Team zum AKH und prescht mit dem Patient zur nicht
vorinformierten Aufnahme.

Der dortige Aufnahmearzt Dr. Wasbringst hat schon einen 22 Stunden
Dienst hinter sich und ist alles andere als erfreut \"uber das
{\quotedblbase}Mitbringsel{\textquotedblleft} des Teams. Wutentbrannt
br\"ullt er Natascha an, was sie sich eigentliche erlaube trotz
gesperrten Bettenkontingent seine Station anzufahren und verweigert ihr
die Aufnahme des Hr. Schmerz.

\paragraph*{Rechtsfrage}

Besteht eine Behandlungspflicht f\"ur Krankenanstalten? Besteht eine
Behandlungspflicht f\"ur medizinische Berufe? Besteht eine
ersthelferische Versorgungspflicht f\"ur den Durchschnittsb\"urger?

\paragraph*{Rechtslage}

\begin{itemize}
\item {\mdseries
	\gEs{Behandlungspflicht f\"ur Krankenanstalten} (=
	KA) - \gZ{{\S}{\S} 22 und 40 KAKuG}:}


	\begin{itemize}
	
		\item {\mdseries
		Jede Krankenanstalt, ob \"offentlich oder privat, ist verpflichtet einen
		behandlungspflichtigen Patienten  aufzunehmen und
		zumindest ärztliche Grundbehandlung zukommen zu lassen. }

		\item {\mdseries
		Das in Wien geltende Bettenkontingentsystem, mit seinen Zuweisungen zu
		bestimmten Krankenanstalten, ist trotzdem zulässig, da
		sich der Patient während des Transportes schon in einem präklinischen Umfeld (=
		Einsatzwagen mit Sanitätern und Medizinprodukten) befindet. Sollte
		sich aber der Zustand des Patienten während der Fahrt
		verschlechtern, darf das nächstgelegene Spital angefahren werden. \newline
		Erst wenn der Patient wieder transportfähig ist, darf er mit geeigneten
		Krankentransportmitteln in eine andere Krankenanstalt weitergeschickt werden.}

	\end{itemize}

	\item {\mdseries
	\gEs{\"Arztliche Behandlungspflicht} --	\gZ{{\S} 48 \"ArzteG}:}

	\begin{itemize}
	
		\item {\mdseries
		\"Arzte, unabhängig ihrer Fachausbildung, sind dazu verpflichtet
		hilfsbed\"urftigen Patienten ärztliche Grundversorgung
		zukommen zu lassen und wenn n\"otig weitere Einsatzkräfte nachzufordern oder den Patient in
		weitere medizinische Behandlung zu schicken.}

	\end{itemize}

	\item {\mdseries
	\gEs{Behandlungspflicht f\"ur Sanitäter} -
	\gZ{{\S} 4 iVm {\S} 8 iVm {\S} 9 Abs. 1 Z 1 SanG}:}

	\begin{itemize}

		\item {\mdseries
		Jeder tätigkeits- oder berufsberechtigter Sanitäter, unabhängig
		von seinem Kompetenzstand, ist zur selbstständigen und
		eigenverantwortlichen Versorgung einer hilfsbed\"urftigen Person
		verpflichtet. Dabei muss er, auch au{\ss}erhalb des Dienstes, zumindest
		qualifizierte Erste Hilfe leisten und wenn n\"otig einen Arzt (jeder
		mit jus practicandi) nachfordern.}

	\end{itemize}

	\item {\mdseries
	\gEs{ersthelferische Versorgungspflicht f\"ur
	Durchschnittsb\"urger} - \gZ{{\S} 95 StGB}:}

	\begin{itemize}
	
		\item {\mdseries
		Jeder B\"urger (und nicht mehr tätigkeitsberechtigter Sanitäter) ist
		zur Setzung von Hilfsleistungen verpflichtet, die seinem Wissens- und
		Kenntnisstand entsprechen. Daher ist jeder mit einem
		F\"uhrerschein-Erste-Hilfe-Kurs zur Setzung von zumutbaren Erste-Hilfe
		Ma{\ss}nahmen verpflichtet.}
		
	\end{itemize}

	\item {\mdseries
	\gEs{Unzumutbarkeit der Behandlung} -
	\gZ{{\S} 95 Abs. 2 StGB}:}

	\begin{itemize}
	
		\item {\mdseries
		F\"ur alle Berufe oder Durchschnittsb\"urger ist die Behandlung dann
		nicht zumutbar, wenn:}
	
	\end{itemize}

	\begin{itemize}
		\item {\mdseries
		man das eigene Leben oder die Gesundheit gefährden w\"urde oder}

		\item {\mdseries
		h\"oherwertige Interessen verletzen m\"usste.}

	\end{itemize}

\end{itemize}

\paragraph*{Rechtsfolge}

\gFazit{
	Auch wenn das Bettenkontingent der Krankenanstalt aufgebraucht ist, muss Dr.
	Wasbringst den Patient aufnehmen und zumindest soweit stabilisieren, dass
	er ohne Gefahren in das eigentliche Zielspital weitertransportiert
	werden kann.}


\subsection{Meldepflichtenrecht}

\subsubsection{Fall 15. Ansteckende Krankheiten}

\paragraph*{Sachverhalt}

Als nun Hr. Schmerz nach langem hin und her auf der Internen aufgenommen
wurde, erholen sich Alex, Sepp und Bernd bei einem kurzen Plausch mit
der NEF-Mannschaft. Die Ruhe hält aber nicht lang an, als auch schon
der nächste Einsatz durchgegeben wird, der RTW und NEF,
zum selben BO schickt. 

Dort angekommen finden sie in einer Lacke von Blut die 25 jährige Fr.
Hilflos liegen, die sich selber (querverlaufend) die Pulsadern
aufgeschnitten hat. Als Sepp sie zu Versorgung beginnt, erzählt Fr.
Hilflos, dass ihr Leben {\quotedblbase}sowieso keinen Sinn mehr
hat{\textquotedblleft}, da sie ihr Freund mit \index{HIV}\gE{HIV} angesteckt hat.

\paragraph*{Rechtsfrage}

Bei welchen Erkrankungen besteht eine Meldepflicht? In welchem
Ausma{\ss} besteht eine solche Meldepflicht?

\paragraph*{Rechtslage}

{\mdseries
\gEs{Meldepflichtige Krankheiten}:}

	\begin{itemize}
	\item {\mdseries
	grundsätzlich:}

	\begin{itemize}
		\item {\mdseries
		leicht \"ubertragbare Krankheiten}

		\item {\mdseries
		mit Gefahr von schweren Gesundheits- oder Todesfolgen}

		\item {\mdseries
		Meldepflicht nur wenn es Gesetz oder Verordnung anordnen}

	\end{itemize}

	\item {\mdseries
	nach \gE{Epidemiegesetz}: }

	\begin{itemize}
		\item {\mdseries
		bei blo{\ss}en Verdacht oder sicherer Diagnose einer im Gesetz genannten
		Erkrankung.}

	\end{itemize}

	\item {\mdseries
	nach \gE{Tuberkulosegesetz}:}

	\begin{itemize}
		\item {\mdseries
		bei blo{\ss}en Verdacht auf eine Erkrankung die mit Sicherheit oder
		Wahrscheinlichkeit durch Tuberkulosebakterien verursacht wurde}

	\end{itemize}
	
	\item {\mdseries
	nach \gE{Geschlechtskrankheitengesetz}:}

	\begin{itemize}
		\item {\mdseries
		nur bei Bef\"urchtung, dass sich die Krankheit	weiterverbreiten w\"urde oder sich der Patient einer
	HB entziehen w\"urde.}

	\end{itemize}
	
	\item {\mdseries
	nach \gE{AIDS-Gesetz}:}

	\begin{itemize}
		\item {\mdseries
		\gE{Anonymisierte} Meldepflicht erst ab: manifester
		Erkrankung mit AIDS (= bei sichtbaren Krankheitszeichen;
		nicht schon bei HIV-Infektion) oder Todesfall mit
		erkennbaren AIDS typischen Krankheiten\footnote{s. Kap. Hygiene, \ref{topic:aids}}}

		\item {\mdseries
		es gibt keine Untersuchungs- oder Behandlungspflicht daher darf kein
		beh\"ordlicher Zwang ausge\"ubt werden}

	\end{itemize}

	\item {\mdseries
	nach \gE{Verordnungen} des Gesundheitsministerium:}

	\begin{itemize}
		\item {\mdseries
		da sich ständig neue \"ubertragbare
		Krankheiten\footnote{in letzter Zeit z.B.
		\gE{SARS}/"'\gE{Vogelgrippe}"´} entwickeln, gibt das
		Gesundheitsministerium in unregelmä{\ss}igen Abständen Verordnungen heraus, die eine Meldepflicht beinhalten.}

	\end{itemize}
\end{itemize}

\gEs{Wer ist zur Meldung verpflichtet?}

Nicht jeder Durchschnittsb\"urger ist zur Meldung verpflichtet, sondern
nur besondere Personengruppen worunter aber alle Gesundheitsberufe
fallen.

\gEs{An wen ist zu melden?}

\begin{itemize}

	\item {\mdseries
	bei AIDS: Gesundheitsministerium}

	\item {\mdseries
	ansonsten: Bezirksverwaltungsbeh\"orde bzw. in Wien an MA 15}

\end{itemize}

\paragraph*{Rechtsfolge}

\gFazit{
	Nun ist zwar unser RTW-Team nicht zur Meldung von AIDS an eine Beh\"orde
	verpflichtet, doch sehr wohl zur Meldung an die weiterbehandelnden
	Gesundheitsberufe aufgrund der schon erwähnten Auskunftspflicht gem.
	\gZ{{\S} 7} SanG.}


\subsubsection{Fall 16. Strafbare Handlungen}

\paragraph*{Sachverhalt}

Die Schnittwunden am Handgelenk sind nicht die einzigen Verletzungen der
Fr. Hilflos. Mehrere Hämatome\footnote{\emph{Hämatom:}
Bluterguss.} \"uberziehen den K\"orper und ein gro{\ss}es
Veilchen ziert das rechte Auge der jungen Patienten. Als Natascha auf diese weiteren Verletzungen
eingeht, beginnt Fr. Hilflos weinend zu berichten, dass ihr Freund sie \"ofters geschlagen hat und er ihr mit
dem Tod gedroht hat, falls sie die Polizei verständigen w\"urde.

\paragraph*{Rechtsfrage}

Besteht eine Pflicht zur Anzeige von strafbaren Handlungen? Wenn ja,
f\"ur wen besteht diese Anzeigepflicht?

\paragraph*{Rechtslage}

{\mdseries
\gEs{Anzeigepflicht}}

Nur bestimmte Personengruppen sind dazu verpflichtet bei Kenntnis \"uber
strafbare Handlungen dies der Polizei oder Staatsanwaltschaft per
Anzeige zu melden.

\begin{itemize}
	\item {\mdseries
	ärztliche Anzeigepflicht - \gZ{{\S} 54 Abs. 4 und 5
	\"ArzteG}:}

	{\mdseries
	nur bei Vorliegen folgender Delikte (Verdacht reicht aus) und
	unabhängig vom Willen des Betroffenen:}

	\begin{itemize}
		\item {\mdseries
		T\"otung bzw. schwere K\"orperverletzung}

		\item {\mdseries
 		Misshandlung, Quälen, Vernachlässigung oder sexueller Missbrauch
		von:}

		\begin{itemize}
			\item {\mdseries
			Minderjährigen oder}

			\item {\mdseries
			Volljährigen, die nicht alleine in der Lage sind ihre Interessen
			wahrzunehmen}

		\end{itemize}
	\end{itemize}
	
	\item {\mdseries
	polizeiliche Anzeigepflicht - \gZ{{\S} 24 StPO}}

	{\mdseries
	f\"ur alle Delikte}

\end{itemize}

{\mdseries
\gEs{Anzeigerecht} - \gZ{{\S} 86 Abs. 1 StPO}:}

Jeder Durchschnittsb\"urger und auch sonstige Gesundheitsberufe sind
nicht verpflichtet zur Anzeige von Strafdelikten. Sie haben aber das
Recht eine strafbare Handlung anzuzeigen. 

Trotz der fehlenden Verpflichtung f\"ur den Sanitäter sollte bei jedem
Einsatz, bei dem der Verdacht einer strafbaren Handlung vorliegt,
zumindest die Polizei nachgefordert werden. Die Nachalarmierung kann
man in weiter Analogie mit der aus \gZ{{\S} 4} SanG
entstammenden Pflicht f\"ur den Sanitäter {\quotedblbase}das Wohl der
Patienten [{\dots}] zu wahren{\textquotedblleft} begr\"unden.

\paragraph*{Rechtsfolge}

\gFazit{
	Natascha muss aufgrund ihrer gesetzlichen Anzeigepflicht die Polizei
	nachfordern, die zur Einbringung einer Anzeige den Sachverhalt
	aufnehmen wird. Dies muss sogar dann geschehen, wenn die Patient selber
	keine Anzeige einbringen m\"ochte.}


\subsection{Stra{\ss}enordnung}

\subsubsection{Fall 17. Blaulichtverwendung}

\paragraph*{Sachverhalt}

Nach Durchf\"uhrung der Behandlung und Anzeige gegen den aggressiven
Freund, begibt sich das RTW-Team ersch\"opft zum St\"utzpunkt zur\"uck.
Hungrig und m\"ude quält sich die RTW-Mannschaft durch das
alltägliche Wiener Verkehrschaos, als Sepp auf einmal meint:
{\quotedblbase} Geh, schalt doch das H\"orndl und die Blitzer ei, damit
dee uns endlie Platz machen und wir einr\"ucken
kännan!{\textquotedblleft}

\paragraph*{Rechtsfrage}

Wann darf das Blaulicht samt Folgetonhorn verwendet werden? Wie darf
bzw. muss man sich bei einer Blaulichtfahrt verhalten?

\paragraph*{Rechtslage}

\gEs{Rechtsquellen}: \gZ{{\S}\,26}~StVO und \gZ{{\S}\,107}~KFG

\subparagraph*{Verwendung von Blaulicht und Folgetonhorn:}

\begin{itemize}
	\item grundsätzlich bei: Gefahr in Verzug
	\item im näheren bei:
	\begin{itemize}
		\item Fahrt zum BO
		\item Haltezeit am BO
		\item Fahrt zum AO
	\end{itemize}
\end{itemize}

\subparagraph*{Verhalten bei Blaulichtfahrt:}

\begin{itemize}
	\item Verkehrsgebote und --beschränkungen gelten nicht
	\item Geschwindigkeitsh\"ochstgrenzen gelten nicht
	\item Ampeln gelten als Stopptafeln
	\item absoluter Vorrang vor Nicht-Einsatzfahrzeugen
	\item es d\"urfen weder Personen noch Sachen gefährdet oder beschädigt
	werden
	\item es darf kein Nicht-Einsatzfahrzeug dem Rettungswagen unmittelbar
	nachfahren
	\item Vertrauensgrundsatz ist ung\"ultig (${\rightarrow}$ man darf bei
	Blaulichtfahrten nicht mehr darauf vertrauen, dass sich die anderen
	Verkehrsteilnehmer richtig verhalten; daher ist defensives
	und voraussehendes Fahren verpflichtend)
\end{itemize}

\subparagraph*{Vorrang der Einsatzfahrzeuge untereinander:}

\begin{enumerate}
	\item Rettungsfahrzeuge
	\item Feuerwehrfahrzeuge
	\item Polizeifahrzeuge
	\item sonstige Einsatzfahrzeuge
\end{enumerate}

\subparagraph*{Missbrauch von Blaulicht und Folgetonhorn:}

{\mdseries
Wer Blaulicht oder Folgetonhorn in Situationen verwendet, bei denen
nicht Gefahr in Verzug vorliegt, macht sich gerichtlich strafbar nach
\gZ{{\S} 1} des {\quotedblbase}Bundesgesetzes gegen
den Missbrauch von Notzeichen{\textquotedblleft}.}

\subparagraph*{\gEs{Gebot bei Blaulichtfahrten}:}
\gFazit{
Nur wer noch gr\"o{\ss}ere Sorgfalt und Vorsicht wahrt als bei normaler Fahrt, verhält sich richtig! 

\gEs{Das beinhaltet mitunter eine \underline{geringere} Fahrgeschwindigkeit als
die h\"ochst zulässige!} }

\subparagraph*{Exkurs: \gEs{Anschnallpflicht} --- \gZ{{\S}\,106}~KFG:}

Der Fahrer eines jeden Fahrzeuges ist dazu verpflichtet zu achten, dass
er und seine Passagiere angeschnallt sind. Es steht aber jedem
Passagier frei zu sagen, dass er unangeschnallt mitfahren m\"ochte,
wodurch eine Schuld beim Fahrer entfällt. F\"ur Kranken- und
Rettungstransport besteht zusätzlich gem. \gZ{{\S}\,106 Abs.\,3~Z\,3} KFG die
Ausnahme, dass ein Passagier (Personal wie Patient) unangeschnallt mitfahren darf, wenn es ansonsten
{\quotedblbase}mit dem Zweck der Fahrt unvereinbar
ist{\textquotedblleft}. Trotzdem sollte der Patient aufgefordert werden
sich anzuschnallen bzw. sollte vom Sanitäter angeschnallt werden.
Verweigert der Patient, muss diese Verweigerung dokumentiert werden (am
besten mittels Revers, den der Patient unterschreibt).

Wird man nun im Zuge eines (sogar unverschuldeten) Verkehrsunfalls
verletzt und war man dabei nicht angeschnallt, verliert man einen Teil
des Schandenersatzanspruchs, weil das Fahren ohne angelegten Gurt als
Mitverschulden gewertet wird.

\paragraph*{Rechtsfolge}

\gFazit{
	W\"urde unser RTW-Team tatsächlich Blaulicht und Folgetonhorn
	einschalten nur um schneller einr\"ucken zu k\"onnen, w\"urden sie sich
	f\"ur die falsche Blaulichtfahrt gerichtlich strafbar machen. Bei einem
	auch  unverschuldeten Unfall m\"ussten unsere drei Sanitäter die
	volle Haftung tragen.}


\vspace{1em}
Der weitere Tag verläuft ruhig mit den \"ublichen Patienten
bis zum \ldots

\vspace{1em}
{\centering
\gEs{{\sffamily\Large Feierabend.}}
\par}

\section*{Abk\"urzungsverzeichnis}
\begin{multicols}{2}
\begin{description}
\item[ABGB]  Allgemeines B\"urgerliches Gesetzbuch

\item[AIDS]  acquired immune deficiency syndrome

\item[AO]  Abgabeort

\item[AR]  Arbeitsrecht

\item[BO]  Berufungsort

\item[DNHG]  Dienstnehmer\-haft\-pflicht\-gesetz

\item[EPH-Gestose] 
\"Odem-Protein\-urie-Hypertonie-Schwanger\-schafts\-erkrankung

\item[HB]   Heilbehandlung 

\item[HIV]  human immundeficiency virus

\item[KA]  Krankenanstalt

\item[KAKuG]  Krankenanstalt und
Kur\-anstaltengesetz

\item[KFG]  Kraftfahrgesetz

%\item[KH]  Krankheit

\item[KTW]  Krankentransportwagen

\item[MA]  Magistratsabteilung

\item[MC]  Myocardinfarkt = Herz\-infarkt

\item[MONA]  Morphin + O\gSub{2} +
Nitro\-lingual-Spray + Aspirin/Aspisol medikament\-\"ose
Herz\-infarkt\-be\-handlung

\item[MPG]  Medizinproduktegesetz

\item[NA]  Notarzt

\item[NaCl]  Natrium-Chlorid

\item[NEF]  Notarzteinsatzfahrzeug

\item[NFS]  Notfallsanitäter

\item[NKA]  Notfallsanitäter mit Kompetenzen der
Arzneimittel

\item[NKI]  Notfallsanitäter mit Kompetenzen der
Intubation

\item[NKV]  Notfallsanitäter mit Kompetenzen des
ven\"osen Zugangs

\item[OGH]  Oberster Gerichtshof

\item[OSR]  Oberster Sanitätsrat

%\item[Pat]  Patient

\item[PatVG]  Patienten\-ver\-f\"ugungs\-gesetz

\item[RF]  Rechtsfolge

\item[RFr]  Rechtsfrage

\item[Rl]  Rechtslage

\item[RQW]  Rissquetschwunde

\item[RS]  Rettungssanitäter

\item[RTW]  Rettungstransportwagen

\item[SanG]  Sanitätergesetz

\item[SchE]  Schadenersatz

\item[SchEA]  Schadenersatzanspruch

\item[StGB]  Strafgesetzbuch

\item[StGG]  Staatsgrundgesetz

\item[StPO]  Strafprozessordnung

\item[StrR]  Strafrecht

\item[StVO]  Stra{\ss}enverkehrsordnung

\item[SV]  Sachverhalt

\item[TB]  Tatbestand

\item[UbG]  Unterbringungsgesetz

\item[VfGH]  Verfassungsgerichtshof

\item[VfR]  Verfassungsrecht

\item[VwGH]  Verwaltungsgerichtshof

\item[VwR]  Verwaltungsrecht

\item[ZDG]  Zivildienstgesetz

\item[ZR]  Zivilrecht

\end{description}

\end{multicols}